\chapter{Referencia Rápida de Common Lisp en LisPy}
\label{app:quickref}

\section*{Operadores Aritméticos}

\begin{center}
\begin{tabular}{ll}
\toprule
\textbf{Operación} & \textbf{Sintaxis} \\
\midrule
Suma & \texttt{(+ a b c ...)} \\
Resta & \texttt{(- a b)} \\
Multiplicación & \texttt{(* a b c ...)} \\
División & \texttt{(/ a b)} \\
División entera & \texttt{(floor a b)}, \texttt{(truncate a b)} \\
Módulo & \texttt{(mod a b)} \\
Exponenciación & \texttt{(expt a b)} \\
Raíz cuadrada & \texttt{(sqrt a)} \\
\bottomrule
\end{tabular}
\end{center}

\section*{Comparación}

\begin{center}
\begin{tabular}{ll}
\toprule
\textbf{Operación} & \textbf{Sintaxis} \\
\midrule
Igualdad numérica & \texttt{(= a b)} \\
Menor que & \texttt{(< a b)} \\
Mayor que & \texttt{(> a b)} \\
Menor o igual & \texttt{(<= a b)} \\
Mayor o igual & \texttt{(>= a b)} \\
Identidad & \texttt{(eq a b)}, \texttt{(eql a b)} \\
Igualdad estructural & \texttt{(equal a b)} \\
\bottomrule
\end{tabular}
\end{center}

\section*{Listas}

\begin{center}
\begin{tabular}{ll}
\toprule
\textbf{Operación} & \textbf{Sintaxis} \\
\midrule
Construir cons & \texttt{(cons a b)} \\
Primer elemento & \texttt{(car list)}, \texttt{(first list)} \\
Resto & \texttt{(cdr list)}, \texttt{(rest list)} \\
N-ésimo & \texttt{(nth n list)} \\
Longitud & \texttt{(length list)} \\
Unir & \texttt{(append list1 list2)} \\
Invertir & \texttt{(reverse list)} \\
\bottomrule
\end{tabular}
\end{center}

\section*{Vectores}

\begin{center}
\begin{tabular}{ll}
\toprule
\textbf{Operación} & \textbf{Sintaxis} \\
\midrule
Crear & \texttt{(make-array n)}, \texttt{(vector a b c)} \\
Acceso & \texttt{(aref vector i)} \\
Asignación & \texttt{(setf (aref vector i) val)} \\
\bottomrule
\end{tabular}
\end{center}

\section*{Hash Tables}

\begin{center}
\begin{tabular}{ll}
\toprule
\textbf{Operación} & \textbf{Sintaxis} \\
\midrule
Crear & \texttt{(make-hash-table :test 'equal)} \\
Acceso & \texttt{(gethash key ht)} \\
Asignación & \texttt{(setf (gethash key ht) val)} \\
Eliminar & \texttt{(remhash key ht)} \\
Limpiar & \texttt{(clrhash ht)} \\
Recorrer & \texttt{(maphash fn ht)} \\
\bottomrule
\end{tabular}
\end{center}

\section*{Control de Flujo}

\begin{center}
\begin{tabular}{ll}
\toprule
\textbf{Constructo} & \textbf{Sintaxis} \\
\midrule
Condicional & \texttt{(if test then else)} \\
Condicional múltiple & \texttt{(cond (test1 forms1) ...)} \\
Iteración simple & \texttt{(dotimes (i n) forms)} \\
Iteración en lista & \texttt{(dolist (x list) forms)} \\
LOOP general & \texttt{(loop for x in list collect x)} \\
\bottomrule
\end{tabular}
\end{center}

\section*{Python Bridge}

\begin{center}
\begin{tabular}{ll}
\toprule
\textbf{Operación} & \textbf{Sintaxis} \\
\midrule
Importar módulo & \texttt{(py:import "module")} \\
Llamar función & \texttt{(py:call fn args...)} \\
Acceder atributo & \texttt{(py:attr obj "name")} \\
Acceder índice & \texttt{(py:getitem obj key)} \\
Ejecutar código & \texttt{(py:run "code")} \\
A Lisp & \texttt{(py:->lisp python-obj)} \\
A Python & \texttt{(py:->python lisp-val)} \\
\bottomrule
\end{tabular}
\end{center}

